\حل{یک سیستم علی است اگر خروجی آن در زمان
\(t_0\)، فقط به ورودی‌های همان لحظه و قبل از آن وابسته باشد. یعنی برای حساب کردن
\(y(t_0)\)، نباید لازم باشد مقدار ورودی در زمان‌های بعد از \(t_0\) را بدانیم.

حالا تابع تبدیل را به صورت زیر می‌نویسیم:
\[
    H(s)=\frac{Y(s)}{X(s)}=\frac{N(s)}{D(s)}
\]
در این کسر، \(N(s)\) صورت و \(D(s)\) مخرج است. درجه‌ی صورت و مخرج را به‌ترتیب با
\[
    m=\deg N(s),\qquad n=\deg D(s)
\]
نشان می‌دهیم.

نکته‌ی مهم این است که در تبدیل لاپلاس، ضرب شدن در \(s\) به معنی مشتق گرفتن در
حوزه‌ی زمان است. یعنی اگر در تابع تبدیل عبارتی مثل \(sX(s)\) داشته باشیم، در
زمان چیزی شبیه \(x'(t)\) ظاهر می‌شود. همچنین اگر \(s^2X(s)\) داشته باشیم، در
زمان چیزی شبیه مشتق دوم ورودی ظاهر می‌شود.

حال اگر درجه‌ی صورت از درجه‌ی مخرج بیشتر باشد، یعنی
    \(m>n\),
بعد از ساده کردن یا تقسیم چندجمله‌ای، در تابع تبدیل جمله‌هایی از جنس \(s\)،
\(s^2\)، و به طور کلی \(s^k\) باقی می‌مانند. بنابراین خروجی مستقیما به مشتق‌های
ورودی وابسته می‌شود؛ مثلا ممکن است برای محاسبه‌ی خروجی به \(x'(t)\) یا
\(x''(t)\) نیاز داشته باشیم.

حالا ببینیم چرا این موضوع با علی بودن مشکل دارد. از تعریف مشتق داریم:
\[
    x'(t_0)=\lim_{\Delta t\to 0}
    \frac{x(t_0+\Delta t)-x(t_0)}{\Delta t}.
\]
در این رابطه، اگر \(\Delta t>0\) باشد، مقدار \(x(t_0+\Delta t)\) مربوط به
آینده‌ی لحظه‌ی \(t_0\) است. پس اگر خروجی در لحظه‌ی \(t_0\) برای محاسبه شدن به
مشتق ورودی نیاز داشته باشد، ممکن است لازم شود مقدار ورودی را کمی بعد از \(t_0\)
بدانیم. این دقیقا برخلاف تعریف سیستم علی است.

بنابراین برای این‌که خروجی به مشتق‌های مستقیم ورودی وابسته نشود، نباید درجه‌ی
صورت از درجه‌ی مخرج بیشتر باشد. پس باید داشته باشیم:
\[
    \deg D(s)\ge \deg N(s).
\]
اگر درجه‌ی صورت و مخرج برابر باشند، خروجی ممکن است به خود ورودی \(x(t)\) در همان
لحظه وابسته باشد، که مشکلی ندارد؛ چون وابستگی به زمان حال هنوز علی محسوب می‌شود.
اما اگر درجه‌ی صورت بزرگ‌تر از درجه‌ی مخرج باشد، خروجی به مشتق ورودی وابسته
می‌شود و این می‌تواند سیستم را ناعلی کند.}
